% AI-written, AI cross-reviewed; no human sign-off; not a claim that RH is proved
% RH Notes blueprint source excerpts for Z1.5.01
% Verbatim excerpts in the notation of the parent blueprint. The labelled editions and source ranges appear below.
% See the unit YAML for the current proof status. Earlier-round passages retain their original wording.
\section*{Z1.5.01. Coefficient representations and the Weyl bound}

% Parent source SHA-256: 1af64922b98ed2310744411b1f18b25cc426c97220646321471836e1a3a409d6
% blueprint/ch_P.tex lines 12487 through 12502; r5:small:Z1.5.01
\subsection{Z1.5.01. A coefficient module of controlled dimension}
\label{r5:small:Z1.5.01}
\textit{Complete.} For totally complex $E$ put $d=|\operatorname{Hom}(E,\mathbb C)|$, $A=n^2d$ and $c=((n+1)^2)^A$. For any strictly decreasing integral tuple $m$ with $|m_i|\le H$, we seek one irreducible algebraic coefficient module for every embedding, of dimension at most $c(H+1)^A$, with infinitesimal character the negative of the multiset $\{(n-1)/2+m_i\}$. Put $\bar k=n-1-k$ and $\nu_k=k+1+H-m_{\bar k}$. The right side is nonnegative. Reversing a strictly decreasing integral sequence increases the entries by at least one at each step, so $\nu$ is weakly decreasing and $\nu_k\le n+2H$. If $n=0$ the tuple and all these inequalities are empty.
\par{\footnotesize \nolinkurl{UpperTheory.V.Af.weyl_complex_of_printed_results} (\nolinkurl{Development/AutomorphicArithmetic/AutomorphicVarieties/ModularityArgument/ThresholdArguments/DenseThresholdCases/TWeylC.lean}\allowbreak:10--44). \nolinkurl{UpperTheory.V.Af.infOf} (\nolinkurl{Development/AutomorphicArithmetic/AutomorphicVarieties/AutomorphicForms/IsoOfFrob.lean}\allowbreak:11). \nolinkurl{UpperTheory.V.Af.WeylC} (\nolinkurl{Development/AutomorphicArithmetic/AutomorphicVarieties/AutomorphicForms/LocIsotypic.lean}\allowbreak:111--113).\par}

For $0\le h<n+2H$ the set $\{k:h<\nu_k\}$ is an initial segment. Multiply the leading minors of these sizes, taking the empty minor to be one. This polynomial is nonzero at the identity. Right multiplication by an upper-triangular matrix $b$ multiplies it by $\prod b_{kk}^{\nu_k}$, since the leading block sees the leading block of $b$. It lies in the stable space of polynomials of degree at most $N=n(n+2H)$, of dimension at most $(N+1)^{n^2}$ by bounding every monomial exponent separately. Its matrix orbit spans a nonzero submodule. Differentiating the triangular character gives diagonal eigenvalues $\nu_k$ and annihilation by the strictly upper generators.
\par{\footnotesize \nolinkurl{UpperTheory.V.Weyl.topPol} (\nolinkurl{Development/AutomorphicArithmetic/AutomorphicVarieties/GaloisLibrary/WeylAlgebraQuantumDet/TopPol.lean}\allowbreak:7--79). \nolinkurl{UpperTheory.V.Weyl.degRank} (\nolinkurl{Development/AutomorphicArithmetic/AutomorphicVarieties/GaloisLibrary/WeylAlgebraQuantumDet/DegRank.lean}\allowbreak:5--30). \nolinkurl{UpperTheory.V.Weyl.existsScalar} (\nolinkurl{Development/AutomorphicArithmetic/AutomorphicVarieties/GaloisLibrary/WeylAlgebraQuantumDet/ExistsScalar.lean}\allowbreak:8--26). \nolinkurl{UpperTheory.V.Weyl.scalar} (\nolinkurl{Development/AutomorphicArithmetic/AutomorphicVarieties/GaloisLibrary/WeylAlgebraQuantumDet/Scalar.lean}\allowbreak:12--45).\par}

On this vector the shifted column determinant has eigenvalue $\prod_k(u-k+\nu_k)$. In a nonidentity permutation the largest moved column $j$ has row $\sigma(j)<j$. Every factor to its right already acts diagonally on the vector, then the factor at $j$ kills it. The determinant commutes with all matrix translates. Indeed, for an invertible change of coordinates $P$ with inverse $Q$, expanding each conjugated shifted generator in the determinant gives two index sums. The first signed sum is the determinant of a row selection of $P$ and vanishes on repeated rows. The second involves the corresponding quantum minor. The $\mathfrak{gl}_n$ commutator with the consecutive shifts gives skew-symmetry under adjacent exchanges, so repetitions vanish and the surviving permutations contribute $\det P\det Q=1$. This gives conjugation invariance, hence commutation with translates by invertible matrices. The entries of the commutator are polynomial and vanish on the invertible locus, so commutation extends to all matrices. The determinant therefore acts by that same scalar on the entire orbit span.
\par{\footnotesize \nolinkurl{Formalism.REAL.quantumDetHighestWeight_holds} (\nolinkurl{Formalism/RealPlaceAnalysis/RealPlaceTaylor/QuantumDetHighestWeightHolds.lean}\allowbreak:68--104). \nolinkurl{GlobalTheory.QuantumMinorAdjacentSwap.qminor_swap_adjacent} (\nolinkurl{Development/AbelianVarietyTheory/Proofs/IwasawaDecomposition/QminorSwapAdjacent.lean}\allowbreak:113--140). \nolinkurl{GlobalTheory.QuantumMinorAdjacentSwap.qdet_conj} (\nolinkurl{Development/AbelianVarietyTheory/Proofs/IwasawaDecomposition/QminorSwapAdjacent.lean}\allowbreak:203--250). \nolinkurl{UpperTheory.V.Weyl.qdetComm} (\nolinkurl{Development/AutomorphicArithmetic/AutomorphicVarieties/GaloisLibrary/WeylAlgebraQuantumDet/QdetComm.lean}\allowbreak:7--29). \nolinkurl{UpperTheory.V.Weyl.commAll} (\nolinkurl{Development/AutomorphicArithmetic/AutomorphicVarieties/GaloisLibrary/WeylAlgebraQuantumDet/CommAll.lean}\allowbreak:6--22). \nolinkurl{UpperTheory.V.Weyl.scalar} (\nolinkurl{Development/AutomorphicArithmetic/AutomorphicVarieties/GaloisLibrary/WeylAlgebraQuantumDet/Scalar.lean}\allowbreak:43--64).\par}

At each complex embedding use the resulting polynomial representation twisted by $\det^{-(n+H)}$ and tensor these factors. The holomorphic derivative at a given embedding sees only its own factor, by the complex-coordinate calculation in Section~\ref{r6:review:Z1.5.01}. Its diagonal is shifted by $-(n+H)$. The scalar polynomial is now $\prod_k(u-(n+H)-k+\nu_k)$. At $u+(n-1)/2$ this has factors $u-(n-1)/2-m_{\bar k}$. Reversal permutes the factors and hence gives the required multiset. A nonzero invariant subspace of least dimension is irreducible, inherits the regular matrix action by restriction, and inherits the determinant scalar by intertwining its differentiated action. Its dimension is bounded by $(N+1)^{n^2d}$, since the tensor product has $d$ factors. Finally $n(n+2H)+1\le(n+1)^2(H+1)$. This also supplies the rank-zero bound.
\par{\footnotesize \nolinkurl{UpperTheory.V.WeylC.fac} (\nolinkurl{Development/AutomorphicArithmetic/AutomorphicVarieties/GaloisLibrary/WeylCharacterRegularity/FacApply.lean}\allowbreak:13--18). \nolinkurl{UpperTheory.V.WeylC.eOf_fac} (\nolinkurl{Development/AutomorphicArithmetic/AutomorphicVarieties/GaloisLibrary/WeylCharacterRegularity/FacApply.lean}\allowbreak:80--97). \nolinkurl{UpperTheory.V.WeylC.exists_irr} (\nolinkurl{Development/AutomorphicArithmetic/AutomorphicVarieties/GaloisLibrary/WeylCharacterRegularity/ExistsTens.lean}\allowbreak:14--29). \nolinkurl{UpperTheory.V.WeylC.exists_tens} (\nolinkurl{Development/AutomorphicArithmetic/AutomorphicVarieties/GaloisLibrary/WeylCharacterRegularity/ExistsTens.lean}\allowbreak:31--51). \nolinkurl{UpperTheory.V.WeylC.algInfChar_iff} (\nolinkurl{Development/AutomorphicArithmetic/AutomorphicVarieties/GaloisLibrary/WeylCharacterRegularity/AlgInfCharIff.lean}\allowbreak:51--53). \nolinkurl{UpperTheory.V.Af.weyl_complex_of_printed_results} (\nolinkurl{Development/AutomorphicArithmetic/AutomorphicVarieties/ModularityArgument/ThresholdArguments/DenseThresholdCases/TWeylC.lean}\allowbreak:32--44).\par}

\textit{[Hypothesis, discharged at the use:]} the published wrapper accepts $\mathrm{Hypotheses}'$; its proof introduces that argument without using a field. The counting theorem passes its own automorphic record when it calls the wrapper. The preceding construction accounts for the weight and dimension clauses.
\par{\footnotesize \nolinkurl{UpperTheory.V.Af.weyl_complex_of_printed_results} (\nolinkurl{Development/AutomorphicArithmetic/AutomorphicVarieties/ModularityArgument/ThresholdArguments/DenseThresholdCases/TWeylC.lean}\allowbreak:10--15). \nolinkurl{UpperTheory.V.Af.card_k_of_printed_results} (\nolinkurl{Development/AutomorphicArithmetic/AutomorphicVarieties/ModularityArgument/ThresholdArguments/DenseThresholdCases/TCardK.lean}\allowbreak:18--24).\par}

% blueprint/ch_P.tex lines 12568 through 12580; r6:review:Z1.5.01
\subsection{Z1.5.01. The shifted minor and the complex derivative}\label{r6:review:Z1.5.01} \emph{Original argument on p.~\pageref{r5:small:Z1.5.01}.} 

The polynomial action used in the Weyl bound has derivatives $E_{ab}$. Differentiating the identity for conjugation by a matrix $y$ gives
$\Xi(y)E_{ab}=\sum_{c,d}y_{ca}(y^{-1})_{bd}E_{cd}\Xi(y)$. Differentiating once more at the identity gives $[E_{ab},E_{cd}]=\delta_{bc}E_{ad}-\delta_{da}E_{cb}$. These derivatives explain the adjacent-column cancellation in the quantum determinant. Put $T_{ab}(x)=E_{ab}+x\delta_{ab}$. In
\[
T_{ad}(x+1)T_{cb}(x)+T_{ab}(x+1)T_{cd}(x)
-T_{cd}(x+1)T_{ab}(x)-T_{cb}(x+1)T_{ad}(x)
\]
the two commutators of the $E$'s cancel the four terms coming from the unit shift. Pair each row permutation with its composition with the adjacent transposition. This identity between the two affected factors, with their left and right products unchanged, reverses the sign of the alternating minor when two adjacent columns are exchanged. Adjacent exchanges generate all permutations; exchanging two identical columns makes the minor its own negative over $\mathbf C$.
\par{\footnotesize \nolinkurl{UpperTheory.V.Weyl.adEq} (\nolinkurl{Development/AutomorphicArithmetic/AutomorphicVarieties/GaloisLibrary/WeylAlgebraQuantumDet/LineDeriv.lean}\allowbreak:25--46). \nolinkurl{UpperTheory.V.Weyl.glRel} (\nolinkurl{Development/AutomorphicArithmetic/AutomorphicVarieties/GaloisLibrary/WeylAlgebraQuantumDet/GlRel.lean}\allowbreak:5--18). \nolinkurl{GlobalTheory.QuantumMinorAdjacentSwap.two_rel}, \nolinkurl{qminor_swap_adjacent}, \nolinkurl{qminor_eq_zero} (\nolinkurl{Development/AbelianVarietyTheory/Proofs/IwasawaDecomposition/QminorSwapAdjacent.lean}\allowbreak:63--101,113--173).\par}

For a tangent direction $X$ in the complex place indexed by $\tau$, the factor indexed by the conjugate embedding sees $\bar X$. The holomorphic derivative uses $(d_X-i\,d_{iX})/2$. It leaves the $\tau$ derivative and cancels the conjugate derivative since conjugation sends $iX$ to $-i\bar X$. Factors at a different place see zero. Total complexity ensures that these are the only coordinate cases. Thus each factor supplies its prescribed scalar polynomial at exactly its own embedding. The construction ignores the automorphic input record; both principal-level counting results pass their respective records to this same construction.
\par{\footnotesize \nolinkurl{UpperTheory.V.WeylC.coord_hol} (\nolinkurl{Development/AutomorphicArithmetic/AutomorphicVarieties/GaloisLibrary/WeylCharacterRegularity/IsRegConj.lean}\allowbreak:69--82). \nolinkurl{UpperTheory.V.WeylC.eOf_fac} (\nolinkurl{Development/AutomorphicArithmetic/AutomorphicVarieties/GaloisLibrary/WeylCharacterRegularity/FacApply.lean}\allowbreak:80--95). \nolinkurl{UpperTheory.V.Af.weyl_complex_of_printed_results} (\nolinkurl{Development/AutomorphicArithmetic/AutomorphicVarieties/ModularityArgument/ThresholdArguments/DenseThresholdCases/TWeylC.lean}\allowbreak:10--44). \nolinkurl{UpperTheory.V.Af.card_k_of_printed_results} (\nolinkurl{Development/AutomorphicArithmetic/AutomorphicVarieties/ModularityArgument/ThresholdArguments/DenseThresholdCases/TCardK.lean}\allowbreak:18--36). \nolinkurl{UpperTheory.V.Br.cardP_of_prints} (\nolinkurl{Development/AutomorphicArithmetic/AutomorphicVarieties/ModularityArgument/ThresholdArguments/DenseThresholdCases/BrCardP.lean}\allowbreak:18--36).\par}

